\sys{} enforces safety at two layers: \emph{program checks} restrict what policy code may execute, and \emph{transition checks} validate each requested resource change against driver-owned state before actuation.
% \sys{} 在两个层面实施安全性：\emph{程序检查}限制策略代码可以执行的内容，\emph{转移检查}在驱动执行前，依据驱动持有的状态验证每一次资源变更请求。

\subsubsection{Driver-owned Transition Validation}
% 驱动持有的转移验证

Policy callbacks never mutate resources directly.
% 策略回调从不直接修改资源。
Typed setters record a callback-local request; after the callback returns, driver code validates the recorded request against the current resource state and only then invokes its native actuator.
% 类型化的 setter 记录回调本地的请求；回调返回后，驱动代码依据当前资源状态验证该请求，然后才调用原生执行器。
An identical repeat request is idempotent, while a conflicting second request is latched and rejected.
% 相同的重复请求是幂等的，而冲突的第二个请求会被锁存并拒绝。

\begin{lstlisting}[language={},xleftmargin=0pt,numbers=none]
record(r, destination, position):
  if not r.attempted: store first raw request
  else if r.pair != (destination, position):
    r.conflict = true

finish_access(root, snapshot, request, action):
  assert PMM list lock held
  result = validate identity, current source,
           generation, conflict, and request range
  if action not in {DEFAULT, BYPASS}: preserve
  else if result == APPLY:
    revalidate; commit one checked native move
  else if result == NO_REQUEST and action == DEFAULT:
    perform native tail move
  else: preserve entry state
\end{lstlisting}

For physical-memory-manager (PMM) reordering, validation checks the owning PMM, root identity, source list, and a generation counter under the existing list lock; only used/unused list heads and tails are valid destinations.
% 对于物理内存管理器（PMM）的重排序，验证在现有列表锁下检查所属 PMM、根标识、来源列表和代际计数器；合法目标仅限 used/unused 列表的头部和尾部。
Fallback is operation-specific: an invalid PMM request preserves entry state, while an invalid prefetch action or region falls back to native prefetch behavior.
% 回退行为因操作而异：非法 PMM 请求保留原有状态，非法的预取动作或区域则回退到原生预取行为。
On RTX~5090, live controls confirm the behavior: 42{,}053 native, 131{,}072 legal-bypass, and 41{,}882 invalid requests each complete with exact native fallback for every invalid request and no data mismatch; scheduler initialization validates timeslice and interleave independently across all 16 cells of the live transition matrix.
% 在 RTX~5090 上，实测控制验证了这一行为：42,053 个原生请求、131,072 个合法 bypass 请求和 41,882 个非法请求均完成，每个非法请求都精确回退到原生行为且无数据不一致；调度器初始化在全部 16 个实测转移矩阵单元中独立验证时间片和交错参数。

\subsubsection{SIMT-aware Verification}
% SIMT 感知验证

Device-side programs additionally pass a GPU verifier that applies PREVAIL~\cite{prevail-verifier} with GPU-specific helper and map models, followed by uniformity analysis and SIMT checks.
% 设备端程序还需通过 GPU 验证器：先以带 GPU 特定 helper 和 map 模型的 PREVAIL~\cite{prevail-verifier} 检查，再进行一致性分析和 SIMT 检查。
Forward dataflow tracks register values, stack bytes, pointer provenance, and map identity using \emph{unknown}, \emph{uniform}, and \emph{varying} tags; a worklist propagates instruction effects until stable, and a final pass checks every reachable instruction.
% 前向数据流以\emph{未知}、\emph{一致}、\emph{可变}标记跟踪寄存器值、栈字节、指针来源和 map 标识；工作队列传播指令效果直至收敛，最后一遍检查每条可达指令。

\begin{lstlisting}[language={},xleftmargin=0pt,numbers=none]
verify_gpu(program, maps):
  reject malformed input or unsupported maps
  require PREVAIL bounds/type/termination checks
  initialize entry state and worklist
  while worklist nonempty:
    pc = pop(worklist)
    out = transfer(program[pc], state[pc])
    join out into successors; enqueue changed states
  require uniform conditional-branch operands
  require uniform atomic target addresses
  require uniform map keys and update flags
  require uniform shared-map values and checked outputs
  reject prohibited helpers or unsafe bridge payloads
\end{lstlisting}

The verifier rejects, for example, an eight-byte stack write at offset $-520$, an unbounded helper-dependent loop, a branch on lane ID, lane-derived map keys and shared-map values, and an atomic whose target comes from a per-thread map; the warp-uniform counterparts of each case are admitted, so atomics on uniform targets remain available to policies.
% 验证器会拒绝诸如偏移 $-520$ 处的 8 字节栈写、依赖 helper 的无界循环、基于 lane ID 的分支、lane 导出的 map 键与共享 map 值，以及目标来自每线程 map 的原子操作；每种情况的 warp 一致版本都会被接受，因此策略仍可使用一致目标上的原子操作。
On RTX~5090, strict-mode pairs on the real counter path admit each 13-instruction callback and record all 32{,}768 invocations, while a 15-instruction lane-varying branch is rejected before insertion; a 15-cell loader matrix further shows strict mode rejects an out-of-stack program without consuming a program ID, whereas warning and default modes admit it.
% 在 RTX~5090 上，真实计数路径上的严格模式用例对接受每个 13 指令回调并记录全部 32,768 次调用，而一个 15 指令的 lane 可变分支在插入前被拒绝；15 单元的加载器矩阵进一步表明，严格模式拒绝越界栈写程序且不消耗程序 ID，而警告和默认模式会接受它。
Verification cost is one-time: admitting the 60-instruction \texttt{kernelretsnoop} tool takes 141~ms and the 13-instruction \texttt{threadhist} takes 12~ms; steady-state throughput with strict verification is statistically indistinguishable from unverified execution for both tools.
% 验证开销是一次性的：接受 60 指令的 \texttt{kernelretsnoop} 工具耗时 141~ms，13 指令的 \texttt{threadhist} 耗时 12~ms；两个工具在严格验证下的稳态吞吐量与未验证执行在统计上不可区分。

\paragraph{Failure classes and TCB}
% 失败类别与 TCB
We distinguish four failure classes: programs the verifier rejects, resource requests that transition validation turns into no-ops, valid policies that do not improve performance, and failures outside policy verification entirely (workload OOM, GPU faults, loader or driver bugs), each diagnosed at its own layer.
% 我们区分四类失败：验证器拒绝的程序、被转移验证变为空操作的资源请求、合法但无效的策略，以及完全处于策略验证之外的失败（工作负载 OOM、GPU 故障、加载器或驱动 bug），各自在所属层面诊断。
The TCB includes Linux's verifier and JIT, \sys{}'s driver hooks and validators, and the native actuators; the verified device path additionally includes PREVAIL, the SIMT analysis, helper and map models, and the GPU compilation toolchain.
% TCB 包括 Linux 验证器与 JIT、\sys{} 的驱动钩子与验证器以及原生执行器；经验证的设备路径还包括 PREVAIL、SIMT 分析、helper 与 map 模型以及 GPU 编译工具链。
Across the agent-driven study of \S\ref{sec:eval}, these checks contained 50 safety-relevant events---24 logic bugs, 18 performance regressions, 2 verifier rejections, 2 GPU faults, and 4 other events---with zero OS kernel panics, and failed policies were detached in seconds without application restarts.
% 在 \S\ref{sec:eval} 的智能体驱动研究中，这些检查控制了 50 起安全相关事件——24 个逻辑错误、18 次性能回退、2 次验证器拒绝、2 起 GPU 故障和 4 起其他事件——OS 内核零崩溃，失败策略在数秒内被摘除，无需重启应用。
